\begin{textsample}{POS Dim 3 – 2022\_c4 – Score 1.00 – 2022\_c4/enade\_2022\_jorn\_23f97d6ee8.txt}  \label{ex:f3_pos_174}
CAPÍTULO 5 PERCEPÇÃO DA PROVA Nas análises feitas neste capítulo, constam as percepções dos concluintes da Área de Jornalismo sobre a prova aplicada no Enade/2022.
Essas percepções foram mensuradas por meio de nove questões que avaliaram desde o grau de dificuldade da prova até o tempo gasto para concluí -la.
As percepções sobre a prova foram relacionadas ao desempenho dos estudantes, à grande região de funcionamento do curso, à categoria administrativa e ao tipo de organização acadêmica da IES.
O questionário de percepção da prova encontra -se ao final do Anexo VIII, no qual está a reprodução do exame.
Na apresentação dos dados relativos às nove questões sobre as percepções a respeito da prova, o desempenho dos estudantes foi classificado em quatro quartos.
Para tanto, esse desempenho foi ordenado de forma ascendente.
O percentil 25, P25, também conhecido como primeiro quartil, é a nota de desempenho que deixa um quarto ( 25 \% ) dos valores observados abaixo e três quartos acima.
O quarto inferior de desempenho é composto pelas notas abaixo do primeiro quartil.
Já o percentil 75, P75, também conhecido como terceiro quartil, é o valor para o qual há três quartos ( 75 \% ) dos dados abaixo e um quarto acima dele.
O quarto superior de desempenho é composto pelas notas iguais ou acima do terceiro quartil.
O percentil 50, P50, também conhecido como mediana, é o valor que divide as notas em dois conjuntos de igual tamanho.
O segundo quarto inclui os valores entre o primeiro quartil ( P25 ) e a mediana.
O terceiro quarto contém os valores entre a mediana ( P50 ) e o terceiro quartil ( P75 ).
Vale ressaltar que percentis, quartis e medianas são pontos que, não obrigatoriamente, pertencem ao conjunto original de dados, ao passo que os quartos são subconjuntos dos dados originais.
Na Figura 5.1, apresenta -se uma ilustração dos quatro quartos descritos.
A seguir, serão apresentados gráficos com resultados selecionados, relativos às nove questões avaliadas por grupos de estudantes.
Nas barras dos gráficos, apresenta -se o percentual de estudantes que assinalaram uma das opções ou a soma das porcentagens daqueles que assinalaram duas ( ou três ) delas.
Por exemplo, para as questões 1 e 2, nos gráficos, é apresentada a porcentagem total de participantes que assinalaram as opções ( D ) Difícil e ( E ) Muito difícil.
Em cada barra, foram assinalados, também, os extremos do intervalo de confiança de 95 \% como linhas verticais unidas por uma linha horizontal na forma da letra H maiúscula.
O estimador de um parâmetro com um certo nível de confiança ( por exemplo, 95,0 \% ) deve conter o parâmetro no intervalo de confiança em 95 \% das vezes.
Na comparação entre os estimadores dos parâmetros de duas classes de uma dada categoria ( por exemplo : Norte e Nordeste nas grandes regiões, de primeiro e último quarto dentro de desempenho, de IES públicas e privadas, ou de universidades e faculdades ), associados aos seus respectivos intervalos de confiança, diz -se que não há diferença estatisticamente significativa entre os parâmetros das duas categorias quando há interseção entre os intervalos de confiança, e que há diferença, se os intervalos de confiança são disjuntos ( para maiores informações vide Glossário ).
Nas Tabelas do Anexo V, são apresentados os valores absolutos e a distribuição percentual⁹ das alternativas válidas das nove questões, segundo o mesmo recorte de desempenho dos estudantes e da grande região de funcionamento do curso, de categoria administrativa e do tipo de organização acadêmica da IES. 5.1 GRAU DE DIFICULDADE DA PROVA Nesta seção são mostradas as percepções dos estudantes quanto ao grau de dificuldade da prova Enade/2022, e que foram analisadas considerando sua vinculação ao Componente de Formação Geral e ao Componente Específico. ⁹ Reitera -se que uma das convenções para tabelas numéricas ( Subseção 2.2.3 do Capítulo 2 ) refere -se à possibilidade de, por questão de arredondamento, a soma das partes não resultar em 100 \%. 5.1.1 COMPONENTE DE FORMAÇÃO GERAL O Gráfico 5.1 apresenta as respostas obtidas para a Questão 1 : “ Qual o grau de dificuldade desta prova na parte de Formação Geral?
De acordo com os dados do Gráfico 5.1, a Região Norte apresentou o maior percentual de estudantes que consideraram a prova difícil ou muito difícil, com 12,7 \%, seguida da região Sul com 11,6 \%.
Na Região Nordeste, a média foi de 9,0 \%, no Sudeste, 8,4 \% e, na Região Centro-Oeste, 8,3 \% ( o menor percentual ).
O Gráfico 5.2 traz mais resultados sobre a questão tratada no gráfico anterior, porém por quartos de desempenho.
Os dados do Gráfico 5.2 mostram que o percentual de estudantes que consideraram o componente de Formação Geral da prova “ Difícil ” ou “ Muito difícil ” foi decrescente em função dos quartos de desempenho : 15,1 \% no 1º quarto e 4,9 \% no 4º quarto, o grupo de melhor desempenho na prova.
Nos quartos de desempenho intermediários, a proporção de estudantes que consideraram a prova difícil ou muito difícil foi de 10,2 \%, no 2º quarto, e de 7,0 \%, no 3º quarto.
O Gráfico 5.3 traz outros resultados também em relação à mesma questão retratada no gráfico anterior, porém por organização acadêmica das IES.
Como mostram os dados do Gráfico 5.3, o percentual de estudantes que consideraram o componente de Formação Geral da prova difícil ou muito difícil foi maior entre os estudantes de Centros universitários, alcançando 11,8 \%, enquanto entre os estudantes das Faculdades alcançou 8,7 \% e em Universidades, a organização acadêmica com a menor incidência, o percentual alcançou 8,6 \%.
O Gráfico 5.4 apresenta os dados por categoria administrativa das IES.
Os dados do Gráfico 5.4 mostram que o percentual de estudantes que consideraram o componente de Formação Geral da prova difícil ou muito difícil foi maior entre os estudantes de cursos de IES privadas ( 10.1 \% ). 5.1.2 COMPONENTE DE CONHECIMENTO ESPECÍFICO O Gráfico 5.5 apresenta as respostas obtidas para a Questão 2 : “ Qual o grau de dificuldade desta prova na parte de Componente Específico? ”.
De acordo com os dados do Gráfico 5.15, a Região Norte apresentou o maior percentual de estudantes que consideraram o componente específico da prova difícil ou muito difícil, com 17,4 \%, seguida da região Sul com 17,0 \%.
Essas regiões tiveram valores máximos de 20,7 \% e 19,4 \%, respectivamente.
Na Região Centro-Oeste a média foi de 14,3 \%.
Nas regiões Sudeste e Nordeste, grupos de menor incidência, as médias foram 12,6 \% e 12,5 \%.
O Gráfico 5.6 traz mais resultados sobre a questão tratada no gráfico anterior, porém por quartos de desempenho.
Os dados do Gráfico 5.6 mostram que o percentual de estudantes que consideraram o componente Específico da prova difícil ou muito difícil foi maior no primeiro quarto de desempenho ( 17,6 \% ) e, no 4º quarto superior ( o grupo de melhor desempenho na prova ), alcançou 10,2 \%.
Nos quartos de desempenho intermediários, a proporção de estudantes que consideraram a prova difícil ou muito difícil foi de 14,0 \% no 2º quarto, e de 12,9 \% no 3º quarto.
O Gráfico 5.7 traz outros resultados também em relação à mesma questão retratada no gráfico anterior, porém por organização acadêmica das IES.
Como mostram os dados do Gráfico 5.7, o percentual de estudantes que consideraram o componente de Formação Geral da prova difícil ou muito difícil foi maior entre os estudantes de Centros universitários, alcançando 15,6 \%, enquanto entre os estudantes das Faculdades, a organização acadêmica com a menor incidência, o percentual alcançou 12,1 \%.
Em Universidades, o percentual intermediário foi 13,3 \%.
O Gráfico 5.8 apresenta os dados por categoria administrativa das IES.
Os dados do Gráfico 5.8 mostram que o percentual de estudantes que consideraram o componente Específico da prova difícil ou muito difícil foi maior entre os estudantes de cursos de IES privadas ( 14,4 \% ).
Enquanto a proporção registrada em IES públicas foi de 12,3 \%. 5.2 EXTENSÃO DA PROVA EM RELAÇÃO AO TEMPO TOTAL O Gráfico 5.9 apresenta as respostas obtidas para a Questão 3 : “ Considerando a extensão da prova, em relação ao tempo total, você considera que a prova foi longa ou muito longa? ”.
Em relação aos dados do Gráfico 5.9, entre as grandes regiões, os que avaliaram a prova como longa ou muito longa, em relação ao tempo total destinado à sua resolução, atingiram uma porcentagem que variou de 35,0 \%, nas regiões Centro-Oeste e Nordeste, até 40,4 \%, na região Sul.
Nas demais regiões, os percentuais foram : 34,9 \% no Norte e 37,6 \% no Sudeste.
O Gráfico 5.10 traz mais resultados sobre a questão tratada no gráfico anterior, porém por quartos de desempenho.
No Gráfico 5.10, observa -se que a proporção de estudantes que consideraram a prova “ longa ” ou “ muito longa ” foi menor no 4º Quarto de desempenho ( 34,3 \% ), e maior, no 1º Quarto inferior ( 41,1 \% ).
Os demais quartos tiveram percentuais : 38,2 \% ( 2º quarto inferior ) e 35,0 \% ( 3º quarto superior ).
O Gráfico 5.11 traz outros resultados também em relação à mesma questão retratada no gráfico anterior, porém por organização acadêmica das IES.
De acordo com os dados do Gráfico 5.11, considerando -se a extensão da prova, o percentual de estudantes que avaliaram a prova como longa ou muito longa, em relação ao tempo total destinado à sua resolução, foi maior entre os estudantes de Centros universitários, alcançando 38,6 \%, enquanto entre os estudantes de Faculdades, a organização acadêmica com a menor incidência, esse índice atingiu 36,3 \%.
Entre os participantes de Universidades, o índice atingiu 36,8 \%.
O Gráfico 5.12 apresenta os dados por categoria administrativa das IES.
Os dados do Gráfico 5.12 mostram que o percentual de estudantes que consideraram a prova longa ou muito longa, em relação ao tempo total destinado à sua resolução, foi maior entre os estudantes de cursos de IES privadas ( 38,2 \% ) do que entre os estudantes de IES públicas ( 35,1 \% ). 5.3 COMPREENSÃO DOS ENUNCIADOS DAS QUESTÕES Também organizados em relação ao Componente de Formação Geral e ao Componente Específico, os dados apresentados nesta seção dizem respeito à forma como os estudantes compreenderam os enunciados de questões da prova, considerando as formulações expressas nesses mesmos enunciados. 5.3.1 COMPONENTE DE FORMAÇÃO GERAL O Gráfico 5.13 apresenta as respostas obtidas para a Questão 4 : “ Os enunciados das questões da prova na parte de Formação Geral estavam claros e objetivos? ”.
Os resultados apresentados no Gráfico 5.13 mostram que, na análise regional, a percentagem de estudantes que avaliaram que todos ou a maioria dos enunciados das questões do Componente de Formação Geral estavam “ claros e objetivos ” variou entre 83,6 \% na região Norte a 92,2 \% na região Sul, menor e maior incidência.
As demais regiões apresentaram os percentuais 87,6 \% no Nordeste, 89,9 \% no Sudeste e 91,9 \% na região Centro-Oeste.
O Gráfico 5.14 traz mais resultados sobre a questão tratada no gráfico anterior, porém por quartos de desempenho.
De acordo com os dados do Gráfico 5.14, segundo o desempenho, observa -se que a proporção dos que emitiram essa opinião cresce, conforme o desempenho aumenta.
No quarto superior, a clareza e a objetividade de todos ou da maioria dos enunciados das questões foram percebidas por 95,6 \%.
Enquanto no quarto inferior essa proporção foi de 79,4 \%.
Os demais quartos apresentaram 92,7 \% ( 3º quarto ) e 89,4 \% ( 2º quarto ).
O Gráfico 5.15 traz outros resultados também em relação à mesma questão retratada no gráfico anterior, porém por organização acadêmica das IES.
Os dados do Gráfico 5.15 mostram que a proporção dos presentes que consideraram que todos ou a maioria dos enunciados das questões do Componente de Formação Geral estavam “ claros e objetivos ” foi maior entre os estudantes de universidades, chegando a 90,2 \%, enquanto entre os centros universitários esse percentual foi de 86,6 \%, a organização acadêmica com a menor incidência.
Entre os estudantes de faculdades, o índice atingiu 88,9 \%.
O Gráfico 5.16 apresenta os dados por categoria administrativa das IES.
Os dados do Gráfico 5.16 mostram que o percentual de estudantes que consideraram que todos ou a maioria dos enunciados das questões do Componente de Formação Geral estavam “ claros e objetivos ” foi maior entre os estudantes de cursos de IES públicas ( 91,6 \% ) quando comparado aos estudantes de IES privadas ( 88,1 \% ). 5.3.2 COMPONENTE DE CONHECIMENTO ESPECÍFICO O Gráfico 5.17 apresenta as respostas obtidas para a Questão 5 : “ Os enunciados das questões da prova na parte de Componente Específico estavam claros e objetivos? ”.
A maioria dos estudantes considerou os enunciados das questões do Componente Específico da prova"claros e objetivos".
Essa avaliação foi expressa por percentuais iguais ou inferiores a 90,7 \% em quatro regiões distintas do país.
A região Norte obteve o menor percentual, com 81,8 \%, enquanto a região Sul registrou o maior percentual, com 90,7 \%.
Nas demais regiões, os percentuais foram de 86,0 \% no Nordeste, 88,9 \% no Sudeste e 89,9 \% no Centro-Oeste.
O Gráfico 5.18 traz mais resultados sobre a questão tratada no gráfico anterior, porém por quartos de desempenho.
Os dados do Gráfico 5.18 mostram que a proporção de estudantes que consideraram os enunciados das questões “ claros e objetivos ” apresenta uma tendência crescente em relação ao aumento de desempenho : mais elevada no quarto superior ( 94,0 \% ), se comparada ao quarto inferior de desempenho ( 80,5 \% ).
O Gráfico 5.19 traz outros resultados também em relação à mesma questão retratada no gráfico anterior, porém por organização acadêmica das IES.
Os dados do Gráfico 5.19 mostram que a proporção dos presentes que consideraram que todos ou a maioria dos enunciados das questões do Componente Específico estavam “ claros e objetivos ” foi maior entre os estudantes de faculdades, atingindo 88,8 \%.
Entre os estudantes de centros universitários, a organização acadêmica com a menor incidência, essa proporção atingiu 86,3 \% e sendo intermediária a proporção alcançada entre os estudantes de universidades ( 88,5 \% ).
O Gráfico 5.20 apresenta os dados por categoria administrativa das IES.
Os dados do Gráfico 5.20 mostram que o percentual de estudantes que consideraram que todos ou a maioria dos enunciados das questões do Componente Específico estavam “ claros e objetivos ” foi maior entre os estudantes de cursos de IES públicas ( 88,8 \% ) que entre os de IES privadas ( 87,7 \% ). 5.4 SUFICIÊNCIA DAS INFORMAÇÕES/INSTRUÇÕES FORNECIDAS O Gráfico 5.21 apresenta as respostas obtidas para a Questão 6 : “ As informações/instruções fornecidas para a resolução das questões foram suficientes para resolvê -las? ”.
De acordo com os dados do Gráfico 5.21 e quanto à distribuição de respondentes pelas grandes regiões, observa -se que a proporção de estudantes que consideraram as informações/instruções fornecidas até excessivas ou suficientes em todas ou na maioria das questões foi sempre superior ou igual a 90,7 \% ( corresponde a menor incidência, na região Norte ), chegando a 95,0 \%, na região Sul ( maior incidência ).
O Gráfico 5.22 traz resultados sobre a questão tratada no gráfico anterior, porém por quartos de desempenho.
De acordo com os dados do Gráfico 5.22, levando -se em conta o desempenho dos participantes, observa -se uma tendência crescente das proporções de participantes que avaliaram as informações/instruções como até excessivas ou suficientes em todas ou na maioria das questões.
O percentual foi mais elevado no quarto superior ( 96,6 \% ).
Já no quarto inferior, a suficiência das informações/instruções declarada como até excessiva, em todas ou na maioria das questões foi percebida por 88,5 \% dos respondentes.
O Gráfico 5.23 traz outros resultados também em relação à mesma questão retratada no gráfico anterior, porém por organização acadêmica das IES.
De acordo com os dados do Gráfico 5.23, considerando -se as informações/instruções fornecidas para a resolução das questões, a proporção que considerou que essas eram até excessivas ou suficientes em todas ou na maioria das questões foi maior entre os estudantes de faculdades totalizando 94,3 \%.
As universidades foram a organização acadêmica em que houve a menor proporção : 93,2 \%.
O percentual dos centros universitários foi intermediário, atingindo 93,3 \% das respostas.
O Gráfico 5.24 apresenta os dados por categoria administrativa das IES.
De acordo com os dados do Gráfico 5.24, o percentual de estudantes que consideraram que as informações/instruções fornecidas para a resolução das questões eram até excessivas ou suficientes em todas ou na maioria das questões foi maior entre os estudantes de cursos de IES públicas, com 94,0 \%, contra 93,0 \%, a proporção registrada entre os de IES privadas. 5.5 DIFICULDADE ENCONTRADA AO RESPONDER À PROVA O Gráfico 5.25 apresenta as respostas obtidas para a Questão 7 : “ Você se deparou com alguma dificuldade ao responder à prova?
Qual? ” Nos Gráficos 5.25 a 5.28, apresentados de forma sequenciada mais adiante, são expostos os percentuais de estudantes que apontaram o “ Desconhecimento do conteúdo ” como “ dificuldade ao responder à prova ”.
O Gráfico 5.25 apresenta as respostas obtidas sobre o “ Desconhecimento do conteúdo ” como a principal “ dificuldade ao responder à prova ”.
De acordo com os dados o Gráfico 5.25, o percentual de inscritos e presentes que apontaram o desconhecimento do conteúdo como “ dificuldade ao responder à prova ” não superou 9,7 \%.
Os percentuais variaram de 9,7 \%, na região Norte, a 6,4 \%, na região Centro-Oeste.
O Gráfico 5.26 traz resultados sobre o ponto tratado no gráfico anterior, porém por quartos de desempenho.
De acordo com os dados do Gráfico 5.26, em relação aos quartos de desempenho, o desconhecimento do conteúdo foi a opção escolhida por 9,3 \% dos estudantes do quarto inferior e por 6,4 \% do quarto superior.
O Gráfico 5.27 traz outros resultados também em relação à mesma questão retratada no gráfico anterior, porém por organização acadêmica das IES.
De acordo com os dados do Gráfico 5.27, na análise por organização acadêmica, o percentual de inscritos e presentes que apontaram o Desconhecimento do conteúdo como “ dificuldade ao responder à prova ” não superou 7,6 \%.
O maior percentual foi nas faculdades com 7,6 \%, seguido pelas universidades com 7,5 \% e centros universitários 7,1 \%.
O Gráfico 5.28 apresenta os dados por categoria administrativa das IES.
De acordo com os dados do Gráfico 5.28, em relação às Categorias Administrativas, o Desconhecimento do conteúdo foi a opção escolhida por 7,5 \% dos estudantes, nas IES públicas e 7,4 \% nas IES privadas. 5.6 CONTEÚDOS DAS QUESTÕES OBJETIVAS DA PROVA O Gráfico 5.29 apresenta as respostas obtidas para a Questão 8 : “ Considerando apenas as questões objetivas da prova, você percebeu que não estudou ainda a maioria desses conteúdos? ”.
Na análise por grande região, a proporção de respondentes que escolheram a opção não estudou ainda a maioria desses conteúdos foi pequena.
Observa -se que, nas regiões Norte ( 2,1 \% ) e Centro-oeste ( 1,9 \% ), as proporções foram ligeiramente maiores.
Em contrapartida, com proporção mais baixa, a região Sul ( 1,2 \% ).
Nas regiões Sudeste e Nordeste, os percentuais foram iguais ( 1,8 \% ), porém diferentes quando observados os valores máximos, sendo de 2,5 \% para a região Nordeste e 2,3 \% para a região Sudeste.
O Gráfico 5.30 traz resultados sobre a questão tratada no gráfico anterior, porém por quartos de desempenho.
De acordo com os dados do Gráfico 5.30, considerando -se separadamente as opiniões dos estudantes dos quatro quartos de desempenho, observa -se que, no primeiro quarto inferior, 4,5 \% ofereceram como resposta que não estudaram ainda a maioria desses conteúdos, sendo 0,4 \% os do quarto superior com a mesma resposta.
Nota -se uma tendência decrescente, ou seja, ao se aproximar do grupo de melhor desempenho na prova, os percentuais caem.
O Gráfico 5.31 traz outros resultados também em relação à mesma questão retratada no gráfico anterior, porém por organização acadêmica das IES.
De acordo com os dados do Gráfico 5.31, na análise por organização acadêmica, a proporção de respondentes, que escolheram a opção não estudou ainda a maioria desses conteúdos, foi pequena.
Observa -se que a proporção nas universidades ( 1,5 \% ) foi menor.
Enquanto nos centros universitários, com 2,5 \% e as faculdades 2,3 \%.
O Gráfico 5.32 apresenta os dados por categoria administrativa das IES.
De acordo com os dados do Gráfico 5.32, em relação às Categorias Administrativas, observa -se que, nas IES privadas, 2,2 \% ofereceram como resposta que não estudou ainda a maioria desses conteúdos, e 1,0 \% dos estudantes de IES públicas apresentaram a mesma resposta. 5.7 TEMPO GASTO PARA CONCLUIR A PROVA O Gráfico 5.33 apresenta as respostas obtidas para a Questão 9 : “ Qual foi o tempo gasto por você para concluir a prova? ”.
Considerando -se as cinco grandes regiões brasileiras, os percentuais dos que utilizaram entre duas e quatro horas para finalizar a prova foram, nas regiões Norte e Nordeste, respectivamente 84,8 \% e 78,8 \%.
Nas regiões Sul e Sudeste, os percentuais de estudantes que dispensaram entre duas e quatro horas para concluir a prova foi de 73,2 \% e 69,2 \%, respectivamente.
A região Centro-Oeste ficou com o percentual de 74,0 \%.
O Gráfico 5.34 traz resultados sobre a questão tratada no gráfico anterior, porém por quartos de desempenho.
De acordo com os dados do Gráfico 5.34, uma vez analisadas as alternativas escolhidas pelos estudantes que se situam nos diferentes quartos de desempenho, observa -se uma tendência crescente : uma proporção menor de participantes no quarto inferior ( 59,3 \% ) declarou ter gasto Entre duas e quatro horas para concluir a prova quando comparada com a proporção de participantes no quarto superior ( 82,8 \% ).
O Gráfico 5.35 traz outros resultados também em relação à mesma questão retratada no gráfico anterior, porém por organização acadêmica das IES.
De acordo com os dados do Gráfico 5.35, na análise por organização acadêmica, os percentuais dos que utilizaram Entre duas e quatro horas para finalizar a prova foram 73,7 \% nas universidades ; 72,8 \% nos centros universitários e o menor percentual nas faculdades 69,5 \%.
O Gráfico 5.36 apresenta os dados por categoria administrativa das IES.
De acordo com os dados do Gráfico 5.36, em relação às Categorias Administrativas, observa -se que, nas IES públicas, 78,4 \% declararam ter gasto Entre duas e quatro horas para concluir a prova e 70,1 \% nas IES privadas com a mesma resposta.
Os resultados do Enade 2022 apresentados neste capítulo contribuem para o entendimento de que a consideração da percepção dos estudantes concluintes sobre a prova do Enade é importante tanto para a compreensão do seu desempenho como para a reflexão sistemática e constante a respeito de determinadas realidades relacionadas a esse mesmo desempenho, no país.
Por outro lado, como todos os elementos do trabalho pedagógico institucional mostram -se diretamente relacionados com o desempenho dos estudantes, também ganha relevância conhecer e analisar os resultados da percepção dos coordenadores de curso sobre a mesma prova.
As informações expressas por estudantes e coordenadores sobre questões diversas associadas aos dois componentes da prova – Formação Geral e Conhecimento Específico – podem agregar valor em ações institucionais voltadas para o aperfeiçoamento do instrumento avaliativo, considerando os vários elementos envolvidos em sua elaboração.
Neste contexto, é relevante que os resultados das percepções de estudantes e coordenadores sobre a prova do Enade ganhem peso como um elemento a mais nas análises da oferta dos cursos nas modalidades presencial e a distância, em nível nacional.
Espera -se, assim, que, em uma visão \textbf{formativa} de avaliação, o conhecimento e a análise desses resultados possam vir a contribuir para a construção de alternativas institucionais que venham a induzir a qualidade da educação superior.

% matched lemmas: formativo
\end{textsample}
